\chapter{Results and Evaluation}\label{ch:ch11}

The preceding chapters described what APEX \emph{is} and how each subsystem was
built. This chapter asks a different, harder question: how good is the delivered
artefact, and how do we know? Evaluating a graduation project fairly means
resisting two temptations --- inflating the result with invented benchmarks, and
deflating it out of false modesty. The stance taken here is that every claim must
either trace to a fact already established in the source (a line count, an index, a
configuration constant, a shipped endpoint) or be labelled explicitly as
illustrative or as a future measurement. Nothing in this chapter is a fabricated
user study or a synthetic benchmark presented as if it were measured.

The evaluation proceeds along four axes. \Cref{sec:res-scale} quantifies the
\emph{scale and scope} of what was actually built. \Cref{sec:res-functional} audits
\emph{functional completeness} against the requirements catalogued in
\Cref{ch:ch03}, giving a full traceability of requirement to delivery status.
\Cref{sec:res-usability} evaluates \emph{usability} against Nielsen's heuristics and
positions the System Usability Scale as the instrument for the empirical study that a
classroom deployment would run next. \Cref{sec:res-performance} characterises
\emph{performance} qualitatively, because APEX's dominant
latency term lives inside an external service and was not instrumented under
controlled load. \Cref{sec:res-timeline} places the work on its calendar, and
\Cref{sec:res-verdict} closes with a candid ``what works / what is limited''
reckoning that ties back to the gaps owned throughout the book.

\section{Scale and Scope of the Delivered System}
\label{sec:res-scale}

The most objective measure of a software project's output is the size and shape of
its codebase. APEX is not a prototype or a proof of concept: it is a working,
two-role, mixed-format examination portal with a real database schema, a real
sandboxed grading pipeline, and a real single-page front end. \Cref{fig:loc}
reports its source size by component.

\begin{figure}[htbp]
\centering
\begin{tikzpicture}
  \begin{axis}[
      ybar,
      width=0.82\textwidth,
      height=6.6cm,
      bar width=52pt,
      every axis plot/.append style={bar shift=0pt},
      xmin=0.4, xmax=2.6,
      ymin=0, ymax=19500,
      ytick={0,4000,8000,12000,16000},
      scaled y ticks=false,
      yticklabel style={font=\scriptsize, /pgf/number format/1000 sep={,}},
      ylabel={Lines of code},
      ylabel style={font=\small},
      xtick={1,2},
      xticklabels={Go backend, TypeScript frontend},
      xticklabel style={font=\small},
      nodes near coords,
      nodes near coords style={font=\scriptsize\bfseries,
        /pgf/number format/1000 sep={,}},
      axis lines=left,
      ymajorgrids=true,
      major grid style={apexlight},
      tick align=outside,
    ]
    \addplot+[draw=apexblue!60!black, fill=apexblue!85,
        nodes near coords style={text=apexblue!55!black,
          font=\scriptsize\bfseries, /pgf/number format/1000 sep={,}}]
      coordinates {(1, 6116)};
    \addplot+[draw=apexgreen!55!black, fill=apexgreen!80,
        nodes near coords style={text=apexgreen!40!black,
          font=\scriptsize\bfseries, /pgf/number format/1000 sep={,}}]
      coordinates {(2, 17423)};
  \end{axis}
\end{tikzpicture}
\caption{Source-code size by language and component.}
\label{fig:loc}
\end{figure}


The backend comprises approximately \textbf{6{,}116} lines of Go across the
feature-scoped \texttt{internal/} packages, and the front end approximately
\textbf{17{,}423} lines of TypeScript and TSX under \texttt{frontend/src}. The
roughly $2.8\!:\!1$ ratio of front-end to back-end code is characteristic of the
architecture rather than incidental: APEX deliberately pushes every security- and
correctness-critical decision to the server --- authentication, the three
authorization gates, grading, and all answer state --- and keeps the client a
presentation and orchestration layer. A compact, auditable Go core carrying the
logic that ``fails an exam'', wrapped in a larger but lower-risk React surface that
renders 47 pages of it, is exactly the distribution one wants for an assessment
system where the browser is treated as untrusted.

Line counts alone under-describe a system, so \cref{fig:counts} summarises APEX
along four structural dimensions that better convey its breadth.

\begin{figure}[htbp]
\centering
\begin{tikzpicture}
  \begin{axis}[
    xbar,
    width=0.86\textwidth,
    height=5.2cm,
    xmin=0, xmax=54,
    bar width=13pt,
    y=15pt,
    enlarge y limits=0.25,
    symbolic y coords={Execution languages,Core entities,Route modules,Frontend pages},
    ytick=data,
    nodes near coords,
    nodes near coords align={horizontal},
    every node near coord/.append style={font=\scriptsize, color=apexblue!80!black},
    tick label style={font=\scriptsize},
    xtick=\empty,
    axis x line*=none,
    axis y line*=left,
    y axis line style={draw=none},
    tick style={draw=none},
  ]
    \addplot[
      draw=apexblue!70!black,
      fill=apexblue,
    ] coordinates {
      (6,Execution languages)
      (15,Core entities)
      (13,Route modules)
      (47,Frontend pages)
    };
  \end{axis}
\end{tikzpicture}
\caption{APEX by the numbers.}
\label{fig:counts}
\end{figure}


Each of these numbers is a verified structural fact, not a marketing figure:

\begin{itemize}[leftmargin=1.4em]
  \item \textbf{15 core entities} --- the GORM models registered in
        \texttt{database.AutoMigrate}, from \texttt{User} and \texttt{Class} through
        \texttt{Exam}, \texttt{Problem}, \texttt{TestCase}, \texttt{ExamAttempt},
        \texttt{Submission}, and \texttt{TestResult} to the supporting
        \texttt{Notification}, \texttt{Announcement}, \texttt{Folder}, and
        \texttt{PasswordResetToken}. These are not toy tables: three of them carry
        unique composite indexes (\texttt{idx\_class\_user},
        \texttt{idx\_user\_exam\_attempt}, \texttt{idx\_exam\_class}) that encode
        genuine business invariants at the storage layer.
  \item \textbf{13 route modules} --- the \texttt{RegisterRoutes} calls wired in
        \texttt{main.go}: auth, class, student, exam, problem, testcase, submission,
        execute, notification, profile, teacher, announcement, and folder. Each is a
        self-contained feature package following the same four-layer request
        lifecycle.
  \item \textbf{47 feature pages} --- the front-end surface a user can actually
        reach, spread across the \texttt{auth}, \texttt{courses}, \texttt{dashboard},
        \texttt{exams}, \texttt{grading}, \texttt{landing}, \texttt{playground},
        \texttt{results}, \texttt{settings}, and \texttt{social} feature directories.
  \item \textbf{6 execution languages} --- Python, JavaScript, TypeScript, Java, C,
        and C++, mapped to Judge0 language IDs in the \texttt{LanguageMap} of
        \texttt{internal/judge0/client.go}; an unknown language is rejected as
        \texttt{unsupported language} rather than silently mis-run.
\end{itemize}

Taken together, these figures answer the ``is it real?'' question quantitatively.
APEX is a mid-sized full-stack application whose complexity is concentrated where it
belongs: fifteen persisted entities, thirteen server modules enforcing a uniform
authorization discipline, and a client broad enough to cover the entire
teacher-and-student workflow.

\section{Functional Completeness Against Requirements}
\label{sec:res-functional}

\Cref{ch:ch03} catalogued the requirements as they were actually realised, with each
functional requirement (FR) traced to an implementing component. The evaluation
question here is the inverse of that traceability: taking the requirement set as the
yardstick, how much of it was delivered, and with what caveats? \Cref{tab:res-fr}
answers this for the twelve teacher, eight student, and seven platform functional
requirements. The status column uses three levels: \emph{Delivered} means the
capability ships and is reachable; \emph{Delivered, with limitation} means the
capability ships but a named gap qualifies it; there are no
\emph{Not delivered} rows, because the requirements chapter was written to record
only functionality that reached the shipped system rather than an aspirational
wish-list.

\begin{table}[htbp]
\centering
\caption{Functional-requirement coverage of the delivered system.}
\label{tab:res-fr}
\small
\begin{tabularx}{\textwidth}{@{}L{1.5cm} L{2.9cm} X@{}}
\toprule
\textbf{Req.} & \textbf{Status} & \textbf{Evidence / caveat} \\
\midrule
FR-T1--T3 & Delivered & Class CRUD, the three-type exam builder, and the
folder-organised question bank all ship as the \texttt{class}, \texttt{problem},
\texttt{testcase}, and \texttt{folder} modules. \\
FR-T4 & Delivered & Draft/publish gating is enforced; publishing requires a
\texttt{start\_time}. This is the requirement whose earlier mishandling caused the
\texttt{is\_draft} incident (\cref{ch:ch05}), now fixed and rule-bound. \\
FR-T5--T6 & Delivered & Multi-class assignment validates class ownership; per-exam
shuffle, results-visibility, and passing-score options persist on the \texttt{Exam}. \\
FR-T7 & Delivered & Coding, MCQ (set-equality), and written grading all run through
\texttt{internal/judge0}; unconfigured MCQ correctly falls to \texttt{pending\_review}
rather than failing a student. \\
FR-T8 & Delivered & The pending-grading queue and \texttt{PUT /submissions/:id/grade}
recompute the attempt aggregate on each manual grade. \\
FR-T9--T11 & Delivered & Announcements fan out notifications; the teacher dashboard,
results explorer, per-class statistics, exam preview, and pending-work grade-gating
all ship. \\
FR-T12 & Delivered & Destructive edits trigger the transactional \texttt{reset\_at}
wipe-and-notify path in \texttt{internal/exam/reset.go}. \\
\addlinespace
FR-S1--S4 & Delivered & Invite-code join, the idempotent one-attempt start,
server-side autosave/resume, and per-attempt shuffle all ship. Autosave payloads are
opaque JSON in \texttt{draft\_answers}, rehydrated on resume. \\
FR-S5 & Delivered & The Monaco editor drives six languages with a
run-against-sample-tests action (\texttt{POST /submissions/run}) that persists no
graded submission. \\
FR-S6 & Delivered, with limitation & Submission closes the attempt and grades
server-side, but grading is a fire-and-forget goroutine with no queue or retry ---
the durable-queue gap owned in \cref{ch:ch07}. \\
FR-S7--S8 & Delivered & Per-test result review, teacher feedback, notifications,
unread counts, and profile/preference management all ship, subject to the teacher's
results-release settings. \\
\addlinespace
FR-P1 & Delivered, with limitation & Email/bcrypt login and 24-hour HS256 JWTs ship;
Google sign-in ships when \texttt{GOOGLE\_CLIENT\_ID} is configured. Limitations: no
rate limiting on auth, a six-character password minimum, and the token in
\texttt{localStorage}. \\
FR-P2--P3 & Delivered & Enumeration-safe password reset (single-use SHA-256 token,
one-hour expiry) and transactional account deletion both ship. \\
FR-P4--P5 & Delivered & The in-process reminder scheduler and the once-only ``Exam
Graded'' notification are deduplicated by database columns, idempotent across
restarts. \\
FR-P6 & Delivered & \texttt{GET /healthz} and the \texttt{--healthcheck} self-probe
support health-checking the shell-less image. \\
FR-P7 & Delivered & Idempotent dev migrations plus a production-only versioned
\texttt{golang-migrate} track (\texttt{SKIP\_AUTOMIGRATE=true}) as sole schema
authority. \\
\bottomrule
\end{tabularx}
\end{table}

The picture is one of \emph{high functional completeness with a small, well-understood
set of qualified rows}. Of the 27 functional requirements, 24 are delivered without
caveat and 3 are delivered with an explicitly named limitation. Crucially, the two
limitations that carry real risk --- fire-and-forget grading (FR-S6) and the
un-hardened authentication surface (FR-P1) --- are not surprises uncovered by this
audit; they are the same gaps owned in the security and architecture chapters and
listed in the project's own README. That consistency is itself an evaluation result:
the system's documented behaviour matches its actual behaviour.

The non-functional requirements (NFR-1 through NFR-20) trace to cross-cutting
mechanisms rather than to endpoints, and their delivery was argued in the chapters
that own them: the fail-fast \texttt{JWT\_SECRET} check and three-gate authorization
(security), the bounded connection pool and indexed hot paths (data model), the
single-artefact deploy and twelve-factor configuration~\cite{twelvefactor,distroless}
(deployment), and the idempotent migration and scheduler discipline
(deployment/DevOps). The one NFR that is deliberately deferred to future empirical
work is NFR-12 (measured satisfaction via SUS), discussed next.

\section{Usability Evaluation}
\label{sec:res-usability}

APEX was not subjected to a formal, statistically powered usability study within the
project timeline; claiming otherwise would overstate the evidence. What can be
reported rigorously is an \emph{expert heuristic evaluation}
against Nielsen's ten usability heuristics~\cite{nielsen-heuristics}, tracing each
relevant heuristic to a concrete, shipped interface mechanism rather than to a
subjective impression. \Cref{tab:res-heuristics} presents that mapping for the
heuristics APEX most directly engages.

\begin{table}[htbp]
\centering
\caption{Heuristic evaluation: Nielsen's heuristics traced to shipped mechanisms.}
\label{tab:res-heuristics}
\small
\begin{tabularx}{\textwidth}{@{}L{4.6cm} X@{}}
\toprule
\textbf{Heuristic} & \textbf{APEX mechanism} \\
\midrule
Visibility of system status &
A submission's lifecycle state (\texttt{pending} $\rightarrow$ \texttt{running}
$\rightarrow$ graded) is surfaced in the UI; live sample-test runs report pass/fail,
output, time, and memory before submit; the autosave indicator shows when work was
last persisted. \\
\addlinespace
User control and freedom &
Server-side autosave with resume: an interrupted attempt continues where it left off
rather than restarting, and a dropped connection or switched device does not cost the
student their work (NFR-11). \\
\addlinespace
Error prevention &
A draft exam cannot be published without a \texttt{start\_time}; a submitted attempt
cannot be resubmitted; the one-attempt-per-\texttt{(student, exam)} unique index makes
a duplicate attempt impossible at the storage layer. \\
\addlinespace
Match between system and the real world &
The vocabulary mirrors the classroom: classes, invite codes, exams, questions,
grading, announcements --- the teacher's mental model, not a database schema. \\
\addlinespace
Consistency and standards &
A single component library (shadcn/ui over Radix primitives) and one design system
across all 47 pages; the same editor (Monaco) for playground and exam. \\
\addlinespace
Recognition rather than recall &
The question bank and folders let a teacher compose exams from previously authored
material instead of re-remembering and re-typing it. \\
\addlinespace
Help users recognise and recover from errors &
Login returns a single generic ``invalid email or password'' (which is also a
security choice); per-test result review shows exactly which hidden tests failed and
the actual versus expected behaviour, so a student can learn from a wrong answer. \\
\bottomrule
\end{tabularx}
\end{table}

Two usability decisions deserve particular emphasis because they resolve tensions
rather than merely satisfying a heuristic. First, output normalisation in the grader
(\texttt{normalizeOutput}) is a usability decision disguised as an engineering one: by
stripping trailing whitespace and normalising line endings symmetrically on the
student's output and the expected output, it prevents the single most demoralising
failure mode of naive auto-graders --- a correct answer failed for an invisible
trailing newline. Second, the exclusion of \texttt{pending\_review} answers from the
attempt aggregate means a student never sees their visible score transiently dip while
a written answer awaits human grading; the number they see is accurate at every moment.

\subsection{Toward an Empirical Study: the SUS Instrument}
\label{sec:res-sus}

Heuristic evaluation finds structural problems but cannot produce a defensible
satisfaction number. The natural next step, and the one NFR-12 commits to, is a
System Usability Scale study~\cite{brooke1996sus}: a ten-item questionnaire whose
alternating positive and negative items are each normalised to a 0--4 contribution,
summed, and multiplied by 2.5 to yield a single 0--100 score, benchmarkable against a
well-established industry mean. The protocol APEX is designed to support is a short
scripted task list exercising both roles --- \emph{build an exam}, \emph{assign it},
\emph{take it}, \emph{grade a written answer}, \emph{review results} --- followed by
the SUS questionnaire, with results reported separately by role because the teacher's
authoring workflow and the student's timed-sitting workflow stress different parts of
the interface.

This study is deliberately \emph{not} reported here with invented numbers. Nielsen's
own guidance is that a handful of participants surfaces the majority of an interface's
usability problems~\cite{nielsen-heuristics}, so even a modest cohort would be a
meaningful first measurement; but a modest cohort drawn from the author's own
institution would also carry a self-selection bias that must be acknowledged when the
study is eventually run. The instrument, the protocol, and the analysis recipe are
specified so that the measurement is reproducible when a real classroom deployment
makes a representative sample available.

\section{Performance Characteristics}
\label{sec:res-performance}

Performance is where candour matters most, because APEX's dominant cost is not in
APEX. This section characterises performance qualitatively and labels clearly what
was measured versus what is reasoned from the design.

\subsection{Grading Latency Is Dominated by Judge0}
\label{sec:res-latency}

The end-to-end latency of grading a coding submission is the sum of three terms: the
Go handler's own work (a database write of the \texttt{Submission} row and, later,
\texttt{TestResult} rows), the HTTP round-trip to Judge0, and --- dominating both ---
Judge0's own compile-and-run time inside the \texttt{isolate}
sandbox~\cite{judge0,isolate}. Because student code is compiled and executed under
per-problem limits (default \texttt{2000}\,ms wall time, \texttt{262144}\,KB memory),
a single test case can legitimately consume up to two seconds of sandbox time, and a
problem with several test cases multiplies that. The Go side, by contrast, adds only
the milliseconds of a local database write and an HTTP call, bounded by a 30-second
client timeout per Judge0 request. The architectural consequence, argued in NFR-6,
is deliberate: exam \emph{submission} does not block on grading at all. Each
submission is graded in a background goroutine, so the HTTP response to
\texttt{POST /student/exams/:id/submit} returns immediately with identifiers, and the
student is never left staring at a spinner whose duration is set by how slow their
own code is. The trade-off is the durable-queue gap already named (\cref{ch:ch07}).

The clearest way to state the performance model is this: \emph{APEX's grading
throughput is Judge0's throughput.} Doubling the sandbox's worker capacity doubles the
number of submissions that can be graded per unit time; the Go server is not the
bottleneck for any classroom-sized load, because it does no CPU-heavy work itself. A
controlled load benchmark against a self-hosted Judge0 --- stepping concurrency and
measuring the latency distribution and the throughput plateau at the worker count ---
is the correct experiment to quantify this, and it is registered as future work
rather than reported here with numbers that were not measured.

\subsection{Image Size, Startup, and API Latency}
\label{sec:res-images}

Several qualitative performance properties follow directly from design choices
documented earlier and can be stated without measurement:

\begin{itemize}[leftmargin=1.4em]
  \item \textbf{Container image size.} The backend is a statically linked
        (\texttt{CGO\_ENABLED=0}) Go binary, stripped with \texttt{-ldflags="-s -w"}
        and built \texttt{-trimpath}, copied onto a
        \texttt{gcr.io/distroless/static-debian12} base that ships essentially only
        the binary and CA certificates --- no shell, no package manager, no libc.
        This is about as small and as narrow-attack-surface as a Go service image
        gets. The one avoidable weight is off-image: a 17\,MB pre-built \texttt{seed}
        binary committed at the repository root, convenient for first boot but
        bloating the clone --- an owned gap, not a runtime cost.
  \item \textbf{Startup.} There is no application-server warm-up: the process reads
        configuration once, opens a bounded connection pool (25 max-open, 10
        max-idle), and begins serving. The reminder scheduler waits a 15-second
        warm-up before its first sweep, so background work never contends with the
        boot path. In Compose, the one-shot \texttt{migrate} service completing
        successfully gates the backend, so the server never accepts a request against
        a stale schema.
  \item \textbf{Non-grading API latency.} Endpoints that do not call Judge0 ---
        starting an attempt, autosaving answers, listing exams, reading a result ---
        are ordinary indexed PostgreSQL reads and writes through GORM, backed by the
        unique composite indexes on the hot paths (NFR-9). These are expected to
        respond in the low tens of milliseconds against a co-located database; the
        precise percentiles are a measurement left for the load study, and the figure
        just given is presented as an \emph{illustrative} design expectation, not a
        recorded result.
\end{itemize}

The overarching performance verdict is that APEX is correctly \emph{sized} for its
stated target --- a classroom of tens of concurrent students, not a MOOC of tens of
thousands --- and that its one heavy operation is deliberately isolated behind an
external service whose capacity the operator controls independently.

\section{Project Timeline}
\label{sec:res-timeline}

APEX was executed over the 2025--2026 graduation year. \Cref{fig:gantt} shows the
schedule as it actually unfolded, with the phases overlapping the way real
engineering does rather than proceeding as a tidy waterfall.

\begin{figure}[htbp]\centering
\resizebox{\textwidth}{!}{%
\begin{ganttchart}[
  hgrid,
  vgrid,
  x unit=1.35cm,
  y unit title=0.55cm,
  y unit chart=0.52cm,
  bar height=0.6,
  title height=1,
  canvas/.append style={fill=apexlight, draw=apexgrey, line width=0.4pt},
  bar/.append style={fill=apexblue, draw=apexblue!60!black, line width=0.4pt, rounded corners=1pt},
  bar label font=\scriptsize,
  title label font=\scriptsize\bfseries,
  title/.append style={fill=apexblue!12, draw=apexgrey},
  bar incomplete/.append style={fill=apexcyan},
  group label font=\scriptsize\bfseries,
  group/.append style={fill=apexgrey, draw=apexgrey}
]{1}{8}
  \gantttitle{2025}{4}
  \gantttitle{2026}{4} \\
  \gantttitlelist{9,10,11,12,1,2,3,4}{1} \\
  \ganttbar[bar/.append style={fill=apexcyan,draw=apexcyan!60!black}]{Requirements \& research}{1}{2} \\
  \ganttbar[bar/.append style={fill=apexcyan,draw=apexcyan!60!black}]{System \& data design}{2}{3} \\
  \ganttbar{Backend core (auth, models, routes)}{3}{5} \\
  \ganttbar{Grading engine (Judge0/isolate)}{4}{6} \\
  \ganttbar{Frontend (SPA, Monaco, exam flows)}{4}{7} \\
  \ganttbar[bar/.append style={fill=apexamber,draw=apexamber!60!black}]{Security review \& hardening}{6}{7} \\
  \ganttbar[bar/.append style={fill=apexgreen,draw=apexgreen!60!black}]{Testing \& CI}{6}{7} \\
  \ganttbar[bar/.append style={fill=apexgreen,draw=apexgreen!60!black}]{Deployment (Docker/compose)}{7}{7} \\
  \ganttbar[bar/.append style={fill=apexgrey,draw=apexgrey!60!black}]{Thesis writing}{5}{8}
\end{ganttchart}%
}
\caption{Graduation-project timeline (Gantt).}\label{fig:gantt}
\end{figure}


Three features of the timeline are worth reading off the chart. First, requirements
and design front-load the schedule but do not finish before implementation begins;
the backend core (auth, models, routes) starts while design is still settling, which
is why the requirements chapter is written as a record of what shipped rather than a
speculative contract. Second, the grading engine and the front end run in parallel
across the middle of the year --- the two largest and most independent workstreams ---
reflecting the front-end-heavy code distribution seen in \cref{fig:loc}. Third, and
most importantly for the project's engineering identity, the \emph{security review and
hardening} phase (in amber) runs \emph{before} deployment, not after. The four
findings --- the IDOR on \texttt{/submissions/run}, the stored-XSS via Markdown, the
weak fallback JWT secret, and the HTML injection in emails --- were all found and
fixed pre-ship, by design. Testing, CI, and Dockerised deployment cluster at the end,
and thesis writing runs as a continuous band across the back half of the year.

\section{What Works and What Is Limited}
\label{sec:res-verdict}

An honest evaluation ends with a plain reckoning. The following two lists are the
distilled verdict; neither hides behind the other.

\subsection*{What works}

\begin{itemize}[leftmargin=1.4em]
  \item \textbf{One pipeline, three formats.} Coding, MCQ, and written questions are
        assessed in a single timed sitting behind one grading path, with a single
        attempt aggregate that recomputes correctly over the exam's full problem list
        and excludes in-flight manual grades from the visible total.
  \item \textbf{Fair auto-grading.} Sandboxed execution under enforced time and
        memory limits, symmetric output normalisation that refuses to fail a correct
        answer on whitespace, and a per-test breakdown that makes every mark
        explainable to the student.
  \item \textbf{Defence in depth that fails closed.} Three independent authorization
        gates (\cref{ch:ch08}) plus a fail-fast secret check that refuses to boot with
        a weak \texttt{JWT\_SECRET}, and a security review conducted before shipping
        rather than after an incident.
  \item \textbf{Resilient exam-taking.} Server-side autosave and single-attempt
        resume mean a crash, a dead battery, or a switched device does not cost a
        student their work --- the student environment is treated as untrusted and
        unreliable, and the server holds the truth.
  \item \textbf{Boring, reproducible deployment.} One Go binary carrying three roles,
        a four-service Compose stack with schema-migration ordering, distroless images
        running as non-root, twelve-factor configuration, and a CI gauntlet that gates
        every change.
  \item \textbf{A schema that learned from its own failure.} The \texttt{is\_draft}
        incident (\cref{ch:ch05}) produced not just a fix but a rule set
        (\texttt{ORM\_RULES.md}) and a production migration discipline in which a
        versioned track, not the ORM, is the sole schema authority.
\end{itemize}

\subsection*{What is limited}

\begin{itemize}[leftmargin=1.4em]
  \item \textbf{Grading has no durable queue.} \texttt{go judge0.Grade(...)} is a bare
        goroutine with no persistence and no retry; a crash mid-grade can strand a
        submission in \texttt{running} and leave an attempt unfinalised
        (\cref{ch:ch07}). Grading is idempotent, so the fix is a durable work queue ---
        the single highest-value piece of future work.
  \item \textbf{The authentication surface is not hardened for public exposure.} No
        rate limiting on auth (bcrypt cost 10 is the only brake on credential
        stuffing), a six-character password minimum, and a JWT in \texttt{localStorage}
        that XSS could lift. Each is a deliberate, documented decision for the thesis
        snapshot, to be hardened before any public deployment.
  \item \textbf{Test coverage is thin.} The backend ships no \texttt{*\_test.go}
        files and the front end a single sanity test; CI already exercises the runners,
        so adding tests requires no pipeline change, but the current automated
        confidence is low relative to the manual and heuristic evaluation.
  \item \textbf{Usability is heuristically, not empirically, evaluated.} The SUS study
        is specified but not yet run on a representative sample; the numbers must come
        from a real deployment, not from the author's institution alone.
  \item \textbf{Performance is reasoned, not benchmarked.} The Judge0-dominated latency
        model is sound but the controlled load experiment --- concurrency sweep,
        latency distribution, throughput plateau --- remains future work.
  \item \textbf{Operational single points.} The reminder scheduler runs in-process on a
        single instance, and the final production deploy hop is still manual; both are
        acceptable at classroom scale and named as such.
\end{itemize}

\section{Summary}
\label{sec:res-summary}

Measured against its own requirements, APEX is a substantially complete system:
roughly 6{,}100 lines of Go and 17{,}400 of TypeScript realise 15 entities, 13 route
modules, and 47 pages, delivering 24 of 27 functional requirements without caveat and
the remaining three with limitations that were owned rather than discovered. Its
usability conforms to the heuristics that matter for an assessment tool, with the
empirical SUS study specified for a future classroom cohort. Its performance model is
bounded by Judge0's sandbox capacity, with the load benchmark registered as
future work rather than fabricated here. The project's engineering identity ---
security before shipping, a schema that learned from a real data-loss incident, and a
deployment that is deliberately boring --- shows up as much in \emph{how} the gaps are
handled as in the features that work. The remaining limitations, each concrete and
each with a known fix, form the agenda carried into \Cref{ch:ch13}.
